\appendix

\section{Proofs of the general results}\label{app:general}

\begin{proof}[Proof of Lemma~\ref{lem:quant}]
If $a\in\bigcap_\lambda E_\lambda(z)$ then each $E_\lambda(z)$ is nonempty, so $z$ lies in the right-hand set. The inclusion fails at $z$ precisely when each $E_\lambda(z)$ is nonempty but $\bigcap_\lambda E_\lambda(z)=\varnothing$, i.e.\ when no common selector exists, which is the stated equality condition.
\end{proof}

\begin{proof}[Proof of Lemma~\ref{lem:cone}]
($\Rightarrow$) $w\ge0$, $w\ne0$, $v\ge0$ give $w\cdot v=\sum_iw_iv_i\ge0$. ($\Leftarrow$) If $v_k<0$, take $w=e_k\in\Wp$ (the $k$-th standard basis vector): $w\cdot v=v_k<0$.
\end{proof}

\begin{proof}[Proof of Proposition~\ref{prop:hierarchy}]
(i) Let $a\in\Etyp(z)$. For every $d$, $\Tube(a,d)\subseteq S=S^{\mathrm{phys}}\cap\{s\ge0\}$; by Lemma~\ref{lem:cone} every tube point satisfies $w\cdot s\ge0$ for every $w\in\Wp$, and likewise every successor lies in $G^w$, so $a\in E_w(z)$. The remaining inclusions follow from $S^w\subseteq S^{\mathrm{phys}}$, $G^w\subseteq G^{\mathrm{phys}}$ and $\End(a,d)\subseteq\Tube(a,d)$.
(ii) The inclusion $\subseteq$ is (i) intersected over $w$. Conversely let $a\in\bigcap_{w\in\Wp}E_w(z)$. For every $d$ and every tube point $p$, $w\cdot s(p)\ge0$ for every $w\in\Wp$, so $s(p)\ge0$ componentwise by Lemma~\ref{lem:cone}; hence $\Tube(a,d)\subseteq S$, and the same argument with $\Succ$ gives $\Succ(a,d)\subseteq G$. Thus $a\in\Etyp(z)$.
\end{proof}

\begin{proof}[Proof of Proposition~\ref{prop:monotone}]
By Proposition~\ref{prop:hierarchy}(i), $E_w(z)\subseteq\Ephys(z)$ for every $w$, so $\mathcal V_w\subseteq\Vphys$ and the intersection over $W$ preserves the inclusion. If $z\in\Vtyp$, Proposition~\ref{prop:hierarchy}(ii) gives $a\in\bigcap_{w\in\Wp}E_w(z)\subseteq\bigcap_{w\in W}E_w(z)$, so $E_w(z)\ne\varnothing$ for every $w\in W$, i.e.\ $z\in\mathcal V_W$. Monotonicity: intersecting over a larger family can only shrink the intersection.
\end{proof}

\begin{proof}[Proof of Theorem~\ref{thm:menu}]
(i) is the margin test restated. For (ii), if $s=d+v$ with $d\in\conv(\mathcal D)$ and $v\ge0$, then for every $w\ge0$, $w\cdot s\ge w\cdot d\ge\min_aw\cdot d_a$, so the scalarized test holds. Conversely suppose $s\notin\conv(\mathcal D)+\Rp$. This set is closed and convex (a compact polytope plus a closed cone), so strict separation gives $w\ne0$ with $w\cdot s<w\cdot d+w\cdot v$ for every $d\in\conv(\mathcal D)$ and every $v\ge0$. Strict separation supplies $w\ne0$. If some coordinate $w_i<0$, choose $v=t e_i$ and let $t\to\infty$; for any fixed $d\in\conv(\mathcal D)$ the right-hand side $w\cdot d+t w_i$ tends to $-\infty$, contradicting the separating inequality. Thus $w\ge0$, so $w\in\Rp\setminus\{0\}$. Taking $v=0$ gives $w\cdot s<\min_{d\in\conv(\mathcal D)}w\cdot d=\min_aw\cdot d_a$, so the scalarized test fails at $w$. The gap formula is (ii) minus (i); if $\mathcal D$ is replaced by $\conv(\mathcal D)$ the union in (i) becomes $\conv(\mathcal D)+\Rp$ and the gap is empty.
\end{proof}

\section{Proofs for the witness datum}\label{app:witness}

\begin{proof}[Proof of Theorem~\ref{thm:witness}]
\textbf{(1)} Typed admissibility requires the worst-case tube in $S_0$ and the successor in $G$. \act{no-switch} never reaches $G$ (its successor retains $q=0$). For \act{fast} the worst-case $s_1$-tube is $[s_1-2,s_1]$, safe iff $s_1\ge2$; its successor $(1,x,s+e)\in G$ given $x\ge0$. Symmetrically \act{slow} is typed-admissible iff $s_2\ge2$. For \act{staged} the floors grow and the successor lies in $G$; the binding constraint is the $x$-tube $[x-1,x]$, safe iff $x\ge1$. Hence $\Vtyp=\{x\ge1\}\cup\{s_1\ge2\}\cup\{s_2\ge2\}$.

\textbf{(2)} Fix $w=(w_1,w_2)\in\Wp$. \act{no-switch} misses $G^w$. \act{fast} is $w$-admissible iff its worst-case aggregate stays nonnegative, $w_1(s_1-2)+w_2s_2\ge0$, i.e.\ $w_1s_1+w_2s_2\ge2w_1$; symmetrically \act{slow} iff $w_1s_1+w_2s_2\ge2w_2$; \act{staged} iff $x\ge1$, since $w\cdot s\ge0$ holds automatically along its tube. Hence $z\in\mathcal V_w$ whenever $x\ge1$; if $x<1$, $z\in\mathcal V_w$ iff one of the two inequalities holds. If $x\ge1$ then $z\in\Vweak$. If $x<1$ and, say, $s_1\ge2$, \act{fast} is typed-admissible by (1), hence $w$-admissible for every $w$, so $z\in\Vweak$. It remains to treat $x<1$, $s_1<2$, $s_2<2$. For $w_1>0$ write $r=w_2/w_1$: \act{fast} is licensed iff $s_1+rs_2\ge2$, i.e.\ $r\ge(2-s_1)/s_2=\rho_1$ (for $s_2>0$; when $s_2=0$ the condition is $s_1\ge2$, excluded), and \act{slow} iff $s_1+rs_2\ge2r$, i.e.\ $r\le s_1/(2-s_2)=\rho_2$. At $w_1=0$ ($r\to\infty$) \act{fast} is licensed provided $s_2\ge0$ while \act{slow} is rejected on $s_2<2$; at $w_2=0$ ($r=0$) symmetrically. Therefore $z\in\mathcal V_w$ for every $w$ iff $\{r\ge\rho_1\}\cup\{r\le\rho_2\}=[0,\infty]$, i.e.\ iff $\rho_2\ge\rho_1$, and
\[
\frac{s_1}{2-s_2}\ge\frac{2-s_1}{s_2}\iff s_1s_2\ge(2-s_1)(2-s_2)\iff s_1+s_2\ge2,
\]
both denominators being positive. Assembling the cases, and noting $\{s_1\ge2\}\cup\{s_2\ge2\}\subseteq\{s_1+s_2\ge2\}$ under nonnegativity, $\Vweak=\{x\ge1\}\cup\{s_1+s_2\ge2\}$.

\textbf{(3)} The physical operators constrain only $x\ge0$ (with $q=1$ at the destination). \act{fast} and \act{slow} have tubes in $S^{\mathrm{phys}}$ and successors in $G^{\mathrm{phys}}$ for every $z\in X_0$ (their dips affect only $s$), so $\Vphys=X_0$; their endpoint values likewise, so $\Vend=X_0$.

\textbf{(4)} By (1)--(2), $\FP=[\{x\ge1\}\cup\{s_1+s_2\ge2\}]\setminus[\{x\ge1\}\cup\{s_1\ge2\}\cup\{s_2\ge2\}]=\{x<1,\ s_1<2,\ s_2<2,\ s_1+s_2\ge2\}=I$, nonempty (e.g.\ $(\tfrac12,\tfrac65,\tfrac65)$) with the stated interior; $R\subseteq\{x\ge1\}\subseteq\Vtyp$.

\textbf{(5)} The first strictness is (4). The point $(\tfrac12,\tfrac1{10},\tfrac1{10})$ has $x<1$ and $s_1+s_2=\tfrac15<2$, so it is outside $\Vweak$ by (2) and inside $\Vphys$ by (3).

\textbf{(6)} The threshold characterisation was derived in (2). On $\Q$, $s_1+s_2\ge2$, so the two licensing intervals overlap and cover all weights; on the relative interior ($s_1+s_2>2$) both thresholds are finite and $r>\rho_2$ licenses \act{fast} only, $r<\rho_1$ \act{slow} only. On $I$ ($x<1$) \act{staged} is inadmissible for every $w$, so $\Etyp(z)=\bigcap_wE_w(z)=\varnothing$ by Proposition~\ref{prop:hierarchy}(ii). On $R$ ($x\ge1$) \act{staged} lies in $\Etyp(z)$ and serves every weight.

\textbf{(7)} On $R$, \act{staged}'s tube is $[x-1,x]\times\{s\to s+e\}$ with $x\ge1$, the floors grow, and the successor $(1,x-1,s+e)\in G$. On $I$ the four violations exhaust the menu: \act{fast}'s $s_1$-tube reaches $s_1-2<0$; \act{slow}'s $s_2$-tube reaches $s_2-2<0$; \act{staged}'s $x$-tube reaches $x-1<0$; \act{no-switch}'s successor misses $G$.
\end{proof}

\begin{proof}[Proof of Theorem~\ref{thm:blend}]
(i) \act{fast}'s worst-case tube has $s_1$-coordinates $[s_1-2,s_1]$ and $s_2$-coordinate $\{s_2\}$; \act{slow}'s is the mirror; both keep $x$ nonnegative. The blended tube therefore has $s_1$-coordinates $[s_1-2\delta,s_1]$, $s_2$-coordinates $[s_2-2(1-\delta),s_2]$, and nonnegative $x$. Typed admissibility requires the tube within $S_0$ and the successor in $G$; the tube condition is exactly $\delta\le s_1/2$ and $\delta\ge1-s_2/2$, and the blended successor is the mixture of two equal successors $(1,x,s+e)\in G$. The window is nonempty iff $1-s_2/2\le s_1/2$, i.e.\ $s_1+s_2\ge2$; on $\Q$ it lies in $(0,1)$ automatically, and at equality it is the singleton $\delta=s_1/2$, feasible because $S_0$ is closed.
(ii) If $s_1+s_2\ge2$, (i) supplies a typed-admissible blend for every $x\ge0$; if $s_1+s_2<2$ no blend is typed-admissible, and the deterministic menu contributes $\Vtyp\subseteq\{x\ge1\}\cup\{s_1+s_2\ge2\}$. Hence the augmented typed set is $\{x\ge1\}\cup\{s_1+s_2\ge2\}=\Vweak$ by Theorem~\ref{thm:witness}(2), and nothing beyond.
(iii) Immediate from the window's dependence on $(s_1,s_2)$ alone and from typed admissibility implying $w$-admissibility for every weight.
\end{proof}

\begin{proof}[Proof of Proposition~\ref{prop:timeshare}]
A policy that spends positive time on \act{fast} visits every $s_1\in[s_1-2,s_1]$, and positive time on \act{slow} every $s_2\in[s_2-2,s_2]$; the union tube lies within $S_0$ exactly when both dips are nonnegative, i.e.\ $s_1\ge2$ and $s_2\ge2$. On $I$ both coordinates are strictly below $2$, so no such policy is typed-admissible.
\end{proof}

\begin{proof}[Proof of Proposition~\ref{prop:rescue}]
Only $\act{staged}_\kappa$ depends on $\kappa$, and its typed admissibility is independent of the floors: its $x$-tube is $[x+\kappa-1,x+\kappa]$, the floors are nondecreasing, and the successor $(1,x+\kappa-1,s+e)\in G$ exactly when $x+\kappa\ge1$. If $z\in\Vtyp$ some base plan is admissible at $\kappa=0$, so $\kappa^*(z)=0$. If $z\notin\Vtyp$ -- $x<1$, $s_1<2$, $s_2<2$ -- then \act{no-switch}, \act{fast} and \act{slow} all fail by Theorem~\ref{thm:witness}(1), the only admissible element of $\mathcal A_\kappa$ is $\act{staged}_\kappa$, admissible exactly for $\kappa\ge1-x$; hence $\kappa^*(z)=1-x$.
\end{proof}

\noindent The proofs of Proposition~\ref{prop:persist}, Theorem~\ref{thm:ladder} and Proposition~\ref{prop:commonshock} are given in SI~S2; the master-equation reduction behind Theorem~\ref{thm:ladder} is the averaging argument that the two principal plans' acceptance conditions are affine in the weight share and sum to a weight-independent quantity, so that the equal weight is always the binding one on the diagonal.
